\chaptertitlepage{Literature Review and Theoretical Framework}{Welfare regimes, subjective late-life experience, and the segmentation tradition}
\label{ch:ch2}

\section{Overview}
\label{sec:ch2-overview}

The analytical strategy of this thesis rests on a single proposition: the over-65 population is not a demographic block but a structured partition of substantively distinct configurations, and the institution that structures that partition is the welfare regime. This chapter builds the proposition from the existing literature and, in doing so, motivates each design choice that the analytical pipeline of Chapter~\ref{ch:ch3} implements. The argument moves from the heterogeneity of the late-life population, to the case for subjective experience as a constituent rather than a derivative dimension of that heterogeneity, to the segmentation tradition that recovers it and the welfare-regime framework that renders the cross-country contrast interpretable; it closes by locating the analytical gap that the empirical chapters address.

\section{Demographic context: population ageing in Europe}
\label{sec:ch2-demographic}

The European demographic transition has produced an over-65 population that is both materially larger and substantially more heterogeneous, on the dimensions of late-life experience, than the cohort of two decades earlier \citep{world2015}. In Italy the over-65 share stands at approximately 24\%, a national over-65 population of 14.18 million in 2024; in Sweden the share is approximately 21\%, or 2.06 million \citep{istat2024, scb2024}. The two countries sit at comparably advanced stages of the transition but under structurally different welfare-state architectures --- the comparative property that the cross-country analysis exploits.

That this heterogeneity is structured, rather than dispersion around a representative trajectory, is the premise on which the segmentation rests. \citet{dannefer2003} formalises the cumulative-disadvantage mechanism, under which working-age inequalities compound across the life course into a more dispersed late-life distribution on resources, health, and subjective dimensions. \citet{netuveli2008} shows that quality-of-life trajectories in the over-50 European population follow distinct latent classes that age-cohort aggregation obscures, and \citet{jivraj2020} demonstrates that the partition into substantively distinct configurations is itself sensitive to the dimensions admitted into the analysis. The heterogeneity is therefore not noise to be averaged away but a partition to be recovered. What structures that partition, the next section argues, is the institutional environment that surrounds the population --- the welfare regime.

\section{Comparative welfare regimes and late-life institutions}
\label{sec:ch2-welfare-regimes}

The comparative welfare-state literature initiated by \citet{espingandersen1990} establishes that the institutional treatment of social risk sorts into a small number of regime types, each producing systematically different outcomes on labour-market participation, family structure, intergenerational transfers, and old-age income. The original three-regime typology --- liberal, conservative-corporatist, and social-democratic --- is built on two properties: de-commodification, the degree to which the welfare state insulates the individual from market dependency, and stratification, the institutional production of class and status differentials.

Two refinements of that typology bear directly on the present analysis. The first is the fourth, Mediterranean-familistic regime articulated by \citet{ferrera1996} and developed by \citet{saraceno2010}, in which the household and the family are the primary loci of social provision in domains the Nordic regime delegates to the public sector. This regime couples a fragmented public-service infrastructure with a residual elderly-care branch and a high level of intra-family responsibility for late-life support; in the cohorts carrying elderly-care obligations, it produces systematically lower female labour-force participation than the Nordic one \citep{bettio2004}. The second refinement is the analytical separation of the social-care branch of the welfare state from its income-transfer branch, drawn by \citet{anttonen1996}: the universal-coverage social-care infrastructure of the Nordic regime --- extensive home- and residential-care delivered by municipal authorities --- is the channel through which the de-familisation of late-life provision operates. This social-care infrastructure is qualitatively distinct from the income transfers on which the Esping-Andersen typology centred.

The empirical signature of the typology is visible across late-life outcomes: old-age poverty risk is more compressed under the Nordic regime than the Mediterranean one \citep{ferrera1996}; healthcare access is more uniformly distributed under universal-access institutions than under fragmented ones \citep{hart1971}; and intergenerational co-residence and informal elderly care are materially more prevalent in the Mediterranean than in the Nordic regime \citep{saraceno2010}. The contrast has a direct operational expression in the formal long-term-care workforce: Sweden employs roughly twelve formal long-term-care workers per 100 people aged 65 or older, against an OECD average of five and an Italian figure near 1.5. The Italian figure reflects the residual care branch of the Mediterranean regime, in which late-life care is discharged predominantly within the family rather than through a municipal workforce \citep{oecd2023}.

This thesis extends the welfare-regime literature at the typological level: the regime contrast operates not only on aggregate outcomes but on the structural shape of the late-life partition that an identical segmentation pipeline returns from each population. The reading developed in Chapter~\ref{ch:ch6} is that the welfare regime is the structuring institution of that partition, and that the structurally different cluster configurations recovered from Italy and Sweden are its empirical signature.

\section{Subjective dimensions of late-life experience}
\label{sec:ch2-subjective-dimensions}

If the welfare regime structures the partition, what dimensions must the partition be drawn on to register that structure? The gerontological literature gives a clear answer: the subjective dimensions of late-life experience carry empirical content that objective state alone does not predict, and they are themselves shaped by the institutional environment. Both propositions are constitutive of the analytical choice made here.

The independence of subjective evaluations is anchored in \citet{idlerbenyamini1997}, whose meta-analysis of twenty-seven cohort studies establishes that self-rated health predicts mortality independently of clinically observed indicators. Reviewing two decades of subsequent work, \citet{bowling2006} concludes that subjective evaluations should be treated as primary rather than derivative late-life indicators. The implication is that segmentations on objective-only inputs systematically discard substantive information. The framing is shared across the active-ageing paradigm of \citet{walker2005active}, the life-course sociology of \citet{lalivedepinay2005}, and the successful-ageing tradition of \citet{rowekahn1997}, each of which treats subjective dimensions as constituent components of the late-life condition rather than as outcomes of objective state.

The measurement instruments carried in SHARE Wave~9 and adopted in the analysis are the operational counterparts of this position. CASP-12 \citep{hyde2003, wiggins2008} measures quality of life; EURO-D \citep{prince1999, castrocosta2008} provides a multi-country comparable measure of depressive affect; and UCLA-3 \citep{hughes2004} measures perceived loneliness on its three-item short form, grounded in the cognitive-discrepancy theory of loneliness \citep{perlman1981}. Each instrument was constructed to capture a substantive component of the late-life condition, not to operationalise an outcome of objective state, with the reliability coefficients and operational detail reported in \S\ref{sec:ch3-subjective-commitment}.

That these dimensions are institutionally shaped is supported by a convergent literature. \citet{steptoe2015} document a social gradient of subjective wellbeing across European older adults, with country-level institutional variation absorbing a substantial share of the cross-country variance; \citet{wahrendorfsiegrist2010} identify welfare-regime variables as the principal cross-country mediators of work-family-late-life trajectories; \citet{courtin2017} document cross-country variation in late-life social isolation that aligns with the welfare-regime typology. The combined evidence licenses the analytical strategy adopted here: the subjective block is a constituent component of the late-life partition, and the cross-country variation observed in the empirical chapters is informative about the regime that surrounds each population.

\section{Multidimensional segmentation of older adults}
\label{sec:ch2-segmentation}

Recovering a partition on a jointly objective-and-subjective input matrix is a segmentation problem, and the methodological literature supplies two traditions. The geometric tradition, initiated by \citet{macqueen1967} and \citet{hartiganwong1979}, partitions a standardised input matrix by minimising the within-cluster sum of squares (K-means); the parametric tradition, developed by \citet{vermuntmagidson2002}, \citet{nylund2007}, and \citet{linzer2011}, recovers the same structure probabilistically under a finite-mixture assumption (Latent Class Analysis). \citet{steinley2006} reviews the properties of K-means, the selection of $k$, and the diagnostics for partition quality; \citet{piccarreta2025} frames the segmentation problem in the demographic and life-course context.

These two traditions are complementary rather than competing: the convergence between a geometric partition and its probabilistic counterpart on the same respondents is itself a validation criterion, high convergence supporting the substantive reading of a partition and low convergence signalling thin cluster boundaries. Upstream of the clustering, factor-analytic preprocessing draws on the classical exploratory-factor-analysis literature --- \citet{kaiser1974} for the rotation criteria and sample-adequacy diagnostics, \citet{horn1965} for the parallel-analysis dimensionality decision. The combined factor-and-segmentation sequence is the standard analytical pipeline in the European-cohort segmentation literature \citep{jivraj2020}.

The decisive design question the literature leaves open is whether the subjective block enters that pipeline as a constituent input or only as external validation. The dominant practice clusters on objective indicators and validates against subjective outcomes \citep{wahrendorfsiegrist2010}; this thesis takes the opposite position, admitting the subjective block as a constituent input. The empirical case for that choice is developed in Chapter~\ref{ch:ch3}.

\section{Silver economy and gerontechnology}
\label{sec:ch2-silver-economy}

The segmentation is not an end in itself: it intervenes in the silver economy, the aggregate of products, services, and institutional arrangements targeted at the over-50 population and the principal commercial-and-policy domain in which the demographic transition of \S\ref{sec:ch2-demographic} becomes an operational problem. The European Commission frames the silver economy as a structural opportunity of the next two decades, with the over-50 cohort generating an estimated 18\% of EU GDP and a market value above EUR 5.7 trillion \citep{vuik2016}. Two structural features of that market bear on the present analysis. The first is the heterogeneity of the addressable population: commercial segmentation that has not internalised the multidimensional structure of late-life experience systematically over-states the addressability of the target cohort. The second is the market's structural dependence on digital adoption --- many silver-economy products that are operationally viable in the Nordic regimes assume an internet-mediated channel that the Mediterranean 65+ cohort does not uniformly support \citep{oecd2023}.

The gerontechnology literature documents the mechanisms behind that dependence. \citet{czaja2006} establish the empirical baseline for adoption barriers in the over-65 population --- perceived self-efficacy, social support for adoption, and prior occupational exposure as the principal predictors. \citet{hargittai2017} shows that the second-level digital divide, the gap in the depth and breadth of online activity conditional on access, is wider among older adults and tracks the same socio-economic gradients that structure late-life resources. \citet{seifert2021} reports a post-pandemic acceleration of digital adoption among European older adults, with country-level heterogeneity that aligns with the welfare-regime contrast of \S\ref{sec:ch2-welfare-regimes}. The digital channel is therefore an empirical variable of the typological partition, not a uniform substrate against which the silver economy can be designed.

The thesis contributes to these literatures on two counts. First, the typologies developed in the empirical chapters translate the multidimensional structure of the over-65 population into addressable segments at a granularity the prevailing commercial reports do not match. Second, the regime-conditioned reading of the digital channel positions the silver-economy intervention in a regime-specific design space rather than a country-aggregate market: the design implications for the under-served Italian segment the analysis identifies follow from the operational meaning of the digital channel under a familistic regime, which the gerontechnology literature documents indirectly but does not typologise.

\section{European typological studies and analytical positioning}
\label{sec:ch2-typologies}

The four strands converge on an analytical position that the existing European typological literature leaves open. Multi-country descriptive analyses on SHARE \citep{jivraj2020} aggregate latent classes across countries and report cross-country share differences, but do not estimate within-country typologies and align them through a structural matching procedure. Country-specific segmentation studies on Italian or Swedish older adults are typically built on objective health and resource indicators alone --- the objective-only specification questioned in \S\ref{sec:ch2-subjective-dimensions} --- and do not extend to the welfare-regime contrast. The protocol adopted here is defined relative to that gap by three explicit choices.

First, the input matrix admits the subjective block on equal footing with the objective indicators, against the more common objective-only clustering with subjective validation; the case rests theoretically on the gerontological tradition reviewed above and is made empirically in Chapter~\ref{ch:ch3}. Second, the analysis is cross-country comparative on a single regime axis --- Mediterranean-familistic versus Nordic-universalistic --- rather than multi-country descriptive: the two-country design trades breadth of coverage for depth of within-comparison interpretation, which the typological extension requires. Third, the comparison is anchored in the welfare-regime literature of \S\ref{sec:ch2-welfare-regimes} rather than in a country-as-control descriptive design, so that the cross-country contrasts are read as institutional signatures rather than residual differences.

These three choices --- subjective inputs as constituent, a single regime axis, a welfare-regime anchor --- together with the structural matching of the two within-country typologies define an analytical instrument that no existing European typological study assembles in full. Table~\ref{tab:ch2-positioning} positions a set of representative studies against the five design dimensions: each existing study realises some of them, none realises all, and the combination is the contribution this thesis develops.

\begin{table}[!htbp]
\centering
\caption[Analytical positioning relative to existing European typological studies]{Analytical positioning of this thesis relative to representative European typological studies, across the five design dimensions that define the analytical instrument. A check mark denotes that the study realises the dimension. No existing study realises all five; the combination is the contribution.}
\label{tab:ch2-positioning}
\small
\setlength{\tabcolsep}{4pt}
\renewcommand{\arraystretch}{1.2}
\begin{tabular}{@{}lccccc@{}}
\toprule
                               & Multidim.   & Subjective  & Cross-      & Welfare-regime & Structural  \\
\textbf{Study}                 & input       & constituent & country     & anchor         & matching    \\
\midrule
Country-specific IT/SE (4D)    & ---         & ---         & ---         & ---            & ---         \\
\citet{netuveli2008}           & \checkmark  & \checkmark  & \checkmark  & ---            & ---         \\
\citet{wahrendorfsiegrist2010} & \checkmark  & ---         & \checkmark  & \checkmark     & ---         \\
\citet{jivraj2020}             & \checkmark  & \checkmark  & \checkmark  & ---            & ---         \\
\textbf{This thesis}           & \checkmark  & \checkmark  & \checkmark  & \checkmark     & \checkmark  \\
\bottomrule
\end{tabular}
\par\medskip
\begin{minipage}{\linewidth}
\footnotesize\textit{Note.} The comparator set comprises representative European typological studies of late-life experience based on SHARE or comparable cohorts, and is illustrative of the design space rather than exhaustive. The check marks reflect the design dimensions realised by each study as reported in its respective publication.
\end{minipage}
\end{table}

\section{Synthesis}
\label{sec:ch2-synthesis}

The argument of this chapter converges on a single analytical instrument: a cross-country, regime-anchored, multidimensional segmentation that admits the subjective block as a constituent input and aligns two within-country typologies through an explicit structural match. Each component is documented in isolation in the existing literature, but no existing study assembles the instrument in full. That assembly is the contribution the next chapter operationalises into the analytical pipeline applied throughout Chapters~\ref{ch:ch4}--\ref{ch:ch7}.
